%%********************Chapter 6**********
\chapter{\uppercase{Result and Conclusion}}

\section{Expected Outcome}
The expected outcome of the proposed project is a mobile-based Financial Goal Planner that provides personalized, adaptive, and risk-aware investment planning for individual financial goals. The system is designed to consider the user's financial goal, target amount, investment horizon, savings capacity, and risk profile while generating suitable investment recommendations.

The proposed system is expected to integrate the following major functionalities:
\begin{itemize}
	\item \textbf{Goal-Based Financial Planning:} The system will allow users to define financial goals by specifying the goal type, target amount, and expected time horizon.
	\item \textbf{Inflation-Adjusted Financial Calculation:} The system will consider the effect of inflation while estimating the future amount required to achieve the selected financial goal.
	\item \textbf{SIP and Savings Calculation:} The system will calculate the required periodic investment and savings contribution based on the target amount, investment horizon, and expected return.
	\item \textbf{Safety and Feasibility Assessment:} The system will evaluate whether the selected financial goal is practically achievable based on the user's available savings, required contribution, investment horizon, and other relevant financial parameters.
	\item \textbf{Risk Profiling:} The system will assess the user's risk tolerance and risk capacity using financial and goal-related information.
	\item \textbf{AI-Driven Recommendation:} The Recommendation Engine will generate personalized investment recommendations by considering the user's financial goal, risk profile, and investment horizon.
	\item \textbf{Dynamic Asset Allocation:} The system will adjust the recommended asset allocation according to the remaining goal horizon. As the goal deadline approaches, the allocation can be planned to become more conservative in order to reduce exposure to higher-risk assets.
	\item \textbf{Mobile Dashboard and Progress Tracking:} The mobile dashboard will provide users with information about their financial goals, required contributions, feasibility status, recommended allocation, and goal progress.
\end{itemize}

The Recommendation Engine will be inspired by the regret-optimized portfolio enhancement and future-looking reward concepts presented in the base paper. These concepts will be adapted to the proposed goal-based financial planning environment rather than directly replicating the original portfolio setting.

The expected system will therefore provide an integrated workflow in which financial calculations, feasibility assessment, risk profiling, investment recommendation, and dynamic allocation are connected through a single mobile-based platform.

\section{Conclusion}
Personal financial planning involves several factors such as financial goals, target amounts, investment horizons, expected returns, inflation, savings capacity, and risk tolerance. Conventional financial planning approaches may use fixed rules or static recommendations and may not adequately consider changes in the remaining goal horizon.

The literature reviewed in this project covers several important areas, including Artificial Intelligence, Reinforcement Learning, Deep Reinforcement Learning, portfolio optimization, robo-advisory, personal financial planning, investor profiling, risk-aware investment, and behavioural finance. These studies demonstrate the potential of intelligent techniques for improving financial decision-making and investment management.

However, the reviewed literature also indicates the need for an integrated approach that connects financial goal planning with feasibility assessment, risk profiling, personalized investment recommendation, and dynamic asset allocation. The proposed project addresses this requirement by designing a goal-oriented financial planning system that combines these components within a single mobile-based platform.

The proposed system is inspired by the base paper, ``Regret-Optimized Portfolio Enhancement through Deep Reinforcement Learning and Future Looking Rewards''. The regret-based and future-looking reward concepts provide the foundation for designing an adaptive recommendation mechanism. In the proposed system, these concepts are considered in the context of individual financial goals, where the investment horizon and risk profile are important factors in determining the recommended allocation.

The major contribution of the proposed system is the integration of multiple financial planning components into one workflow. The system is planned to accept user-specific financial and goal-related inputs, perform the required calculations, assess feasibility, determine the risk profile, generate an investment recommendation, and dynamically adjust the allocation according to the remaining goal horizon.

Phase~1 of the project established the necessary foundation for the subsequent implementation. It included the study of the base paper, detailed literature review, identification of the research gap, formulation of the problem statement, definition of objectives, analysis of system requirements, design of the proposed architecture, identification of major modules, and preparation of the development workflow.

The subsequent phases will focus on the implementation and integration of the proposed modules. These phases will include mobile application development, financial calculation implementation, safety and feasibility assessment, risk profiling, recommendation-engine development, dynamic allocation, backend and database integration, model training and evaluation, system testing, and performance evaluation.

Thus, Phase~1 provides a clear research and technical foundation for developing the proposed Financial Goal Planner as an AI-driven, goal-based, and risk-aware investment planning system.
